% Generated draft: requires mathematical and source review.
\documentclass[11pt]{article}
\usepackage[margin=25mm]{geometry}
\usepackage{fontspec}
\setmainfont{Latin Modern Roman}
\usepackage{amsmath,amssymb,mathtools,unicode-math}
\setmathfont{Latin Modern Math}
\usepackage{enumitem,needspace,xurl,hyperref}
\setlength{\parindent}{0pt}
\setlength{\parskip}{5pt}
\setlength{\emergencystretch}{3em}
\newcommand{\N}{\mathbb{N}}
\newcommand{\Z}{\mathbb{Z}}
\newcommand{\Q}{\mathbb{Q}}
\newcommand{\R}{\mathbb{R}}
\newcommand{\C}{\mathbb{C}}
\newcommand{\F}{\mathbb{F}}
\newcommand{\A}{\mathbb{A}}
\newcommand{\PP}{\mathbb{P}}
\newcommand{\E}{\mathbb{E}}
\newcommand{\eps}{\varepsilon}
\newcommand{\dd}{\,\mathrm{d}}
\begin{document}
\section*{UnsolvedMath ID: 30006034}
\textbf{Source problem code:} catalogue entry
\subsection{Definitions and mathematical statement}
Let \(C_2=\Z/2\Z\). For a genuine \(C_2\)-equivariant spectrum \(X\), write \(\Phi^{C_2}X\) for its geometric fixed-point spectrum. Write \(\mathrm{THR}\) for real topological Hochschild homology. The supplied passage concerns a ring spectrum \(R\) with additional equivariant or duality structure, but does not specify that structure or the hypotheses on \(R\).

The passage reports the equivalence
\[
\Phi^{C_2}\mathrm{THR}(R)
\simeq
\Phi^{C_2}R\otimes_R\Phi^{C_2}R.
\]
Here \(\otimes_R\) denotes a derived relative tensor product, provided the required right and left \(R\)-module structures have been specified. Those structures, the referenced \(C_2\)-action, and the hypotheses under which the displayed equivalence holds are absent from the supplied excerpt. Accordingly, this formula is retained as a reported comparison, not asserted for arbitrary ring spectra.

\textbf{Problem.} Determine how the geometric fixed-point construction and the displayed comparison should extend from the ring-spectrum setting to Poincar\'e \(\infty\)-categories. In particular, formulate the appropriate categorical inputs, the geometric fixed-point object associated with their real topological Hochschild homology, and a categorical counterpart of the right-hand side that specializes to the reported formula under the appropriate ring-spectrum comparison.

A Poincar\'e \(\infty\)-category here means an \(\infty\)-category equipped with Poincar\'e structure in the sense intended by the cited report. The supplied record does not specify its axioms, admissible size conditions, or the allowed class of Poincar\'e structures. It also does not specify the categorical tensor construction or a ring-to-category comparison. Thus the requested extension is a formulation and construction problem; the evidence does not determine a fully quantified conjectural equivalence.
\subsection{Short English statement}
Extend the reported geometric fixed-point formula for real topological Hochschild homology of ring spectra to Poincar\'e \(\infty\)-categories. The source leaves the required structures, hypotheses, and categorical analogue of the tensor product unspecified, so these remain part of the problem.
\subsection{Sources}
\begin{itemize}
\item Supplied local UnsolvedMath registry record, permanent ID \(30006034\), exact source problem code \texttt{catalogue entry}, title: \emph{Geometric Fixed Points for Poincar\'e Infinity-Categories}. The record supplies proposed year 2024 and names no proposer.
\item Supplied report citation: \emph{Homotopical Algebra and Higher Structures (2025)}, Oberwolfach Reports, pp.\ 2251--2308, DOI \texttt{10.4171/owr/2024/39}. \url{https://doi.org/10.4171/owr/2024/39}. The literature-source label is \emph{Oberwolfach Report 39 (2024)}. The supplied passage attributes the displayed ring-spectrum formula to reference [2], whose bibliographic identity is not supplied.
\item Supplied literature source: \emph{Rasmusen, THR of Poincare infinity-categories}. \url{https://younghomseminar.github.io/posts/2023-04-13/}.
\item Literature assessment dated exactly \texttt{2026-08-22}: the supplied curation labels the problem open, describes the categorical extension as a natural question in the report, and states that no proof of the extracted extension was found. This is supplied dated curation derived from the OWR material, not an independently verified literature assessment. Its time-sensitive status and cited claims require independent verification.
\end{itemize}
\end{document}
